\subsection{The ordinary eighth moment of the unit short-one connector}\label{proof:Q5-C049}
On any nonempty finite probability space \((X,\allowbreak \mu)\), let \(K\) be real and antisymmetric with \(|K(x,\allowbreak y)|\le\allowbreak 1\). Delete zero-mass atoms and use the same \(\mu\) for all operators and inner products. Set
\[
\begin{gathered}
Jv=(\E v)\one,\quad A=J+K,\quad T=J-K=A^*,\quad\\
g=K\one,\quad p=1+g,\quad q=1-g,\quad r=Kg,\quad s=K^2g,
\end{gathered}
\]
\[
d=\E g^2,\qquad \delta_2=\E r^2,\qquad
\delta_3=\E s^2,\qquad \theta=13/32.
\]
The finite spectral input in Section~\ref{proof:Q5-C081} gives
\[
0\le d\le1/3,\qquad
\delta_2\le\theta d,\qquad \delta_3\le\theta\delta_2\le\theta^2d.
\tag{Q49.1.1}\label{eq:Q49:1.1}
\]
The actual terminal and unit-connector messages are
\[
f_2=A^2\one=p+r,\qquad F=qf_2,\qquad H=pTF=pT(qf_2).
\tag{Q49.1.2}\label{eq:Q49:1.2}
\]
The terminal unit arm points inward, its two-edge arm points outward, the connector points from the terminal to the receiving root, and the receiving unit arm points outward. These directions are fixed by the tree.
We prove, for every such host,
\[
\E H^8\le1-d/500.
\tag{Q49.1.3}\label{eq:Q49:1.3}
\]
Complete reversal gives
\[
\E\bigl[\{qA(p\bar f_2)\}^8\bigr]\le1-d/500,\qquad
\bar f_2=T^2\one.
\tag{Q49.1.4}\label{eq:Q49:1.4}
\]
Genuine equal pairs can be restored at their roots using \(f_\ell h_\ell\le\allowbreak 1\) and positivity.

\paragraph{Q49.2 A row bound for $K^2$}
For every real \(v\) with \(|v|\le\allowbreak 1\),
\[
\|Kv\|_2^2\le1/3.
\tag{Q49.2.1}\label{eq:Q49:2.1}
\]
To see this, for a legal antisymmetric kernel \(L\), put
\(a=\allowbreak L(x,\allowbreak y)\), \(b=\allowbreak L(y,\allowbreak z)\), \(c=\allowbreak L(z,\allowbreak x)\).
Pointwise,
\[
1+ab+bc+ca
=\frac{(1+a)(1+b)(1+c)+(1-a)(1-b)(1-c)}2\ge0.
\]
Integrating over three independent labels with measure \(\mu\), each quadratic term equals \(-\E(L\one)^2\). Hence
\[
\E(L\one)^2\le1/3.
\tag{Q49.2.2}\label{eq:Q49:2.2}
\]
For \(v\in\{-1,\allowbreak 1\}\), the kernel \(L(x,\allowbreak y)=\allowbreak v(x)K(x,\allowbreak y)v(y)\) is legal and
\(L\one(x)=\allowbreak v(x)Kv(x)\), proving \eqref{eq:Q49:2.1}.
For general \(|v|\le\allowbreak 1\), choose independent signs \(\sigma(x)\) at the finitely many atoms, with \(\E_\sigma\sigma(x)=\allowbreak v(x)\). Convexity gives
\[
\|Kv\|_2^2
=\|\E_\sigma K\sigma\|_2^2
\le\E_\sigma\|K\sigma\|_2^2\le1/3.
\]

For a fixed \(x\), let \(k_x(y)=\allowbreak K(x,\allowbreak y)\). Since \(|k_x|\le\allowbreak 1\), for every real \(w\),
\[
K^2w(x)=\ip{k_x}{Kw}=-\ip{Kk_x}{w},\qquad
|K^2w(x)|^2\le\|Kk_x\|_2^2\|w\|_2^2\le\|w\|_2^2/3.
\tag{Q49.2.3}\label{eq:Q49:2.3}
\]
In particular,
\[
|s(x)|^2\le d/3.
\tag{Q49.2.4}\label{eq:Q49:2.4}
\]
Equation \eqref{eq:Q49:2.1} is an \(L^\infty\)-to-\(L^2\) estimate, while
\eqref{eq:Q49:2.3} is a row estimate for \(K^2\). The spectral bound used below remains \(\theta=\allowbreak 13/32\).

\paragraph{Q49.3 A positive square envelope}
Since \(q,\allowbreak f_2\ge\allowbreak 0\) and \(q+f_2=\allowbreak 2+r\), AM--GM gives
\(F=\allowbreak qf_2\le\allowbreak (1+r/2)^2=\allowbreak 1+r+r^2/4\).
Using \(\E r=\allowbreak -d\), \(Kr=\allowbreak s\), and \(0\le\allowbreak  T(r^2)\le\allowbreak 2\E r^2=\allowbreak 2\delta_2\),
\[
TF\le q-d-s+\delta_2/2\le q-(51/64)d-s.
\tag{Q49.3.1}\label{eq:Q49:3.1}
\]
The last inequality uses \(\delta_2\le\allowbreak 13d/32\), and its right side is nonnegative because it dominates \(TF\).
Put \(a=\allowbreak 51/128\) and \(c=\allowbreak 1-ad\).
AM--GM for \(p\) and the nonnegative upper in \eqref{eq:Q49:3.1} gives
\[
0\le H=pTF\le[1-(51/128)d-s/2]^2=(c-s/2)^2.
\tag{Q49.3.2}\label{eq:Q49:3.2}
\]
Moreover,
\[
c\ge111/128>0,\qquad |s|/2\le\sqrt{d/12}\le1/6<c.
\tag{Q49.3.3}\label{eq:Q49:3.3}
\]
Skewness gives \(\E s=\allowbreak \ip{\one}{K^3\one}=\allowbreak 0\), and
\[
\E s^2=\delta_3\le\theta^2d.
\tag{Q49.3.4}\label{eq:Q49:3.4}
\]

\paragraph{Q49.4 The scalar sixteenth-moment estimate}
At \(d=\allowbreak 0\), \(g=\allowbreak r=\allowbreak s=\allowbreak 0\) and the actual messages \(p,\allowbreak q,\allowbreak f_2,\allowbreak F,\allowbreak H\) are one.
For \(d>\allowbreak 0\), set
\[
x=\frac{\sqrt{d/12}}c,\qquad \eta=-\frac{s}{2c},\qquad
\alpha=3\theta^2=\frac{507}{1024},\qquad
k=12a=\frac{153}{32},\qquad w=\frac1c.
\]
The function \(\sqrt d/(1-ad)\) increases, so
\[
0<x\le\frac{1/6}{111/128}=\frac{64}{333}<\frac15.
\tag{Q49.4.1}\label{eq:Q49:4.1}
\]
The row and second-moment bounds give
\[
|\eta|\le x,\qquad \E\eta=0,\qquad \E\eta^2\le\alpha x^2.
\tag{Q49.4.2}\label{eq:Q49:4.2}
\]
For every integer \(j\ge\allowbreak 2\), \(\eta^j\le\allowbreak  x^{j-2}\eta^2\):
for even \(j\) use absolute values; for odd \(j\) and negative \(\eta\), the left side is nonpositive.
The binomial expansion therefore gives
\[
\E(1+\eta)^{16}
\le1+\alpha[(1+x)^{16}-1-16x].
\tag{Q49.4.3}\label{eq:Q49:4.3}
\]
The exact relation \(1-c=\allowbreak ad\) implies
\[
w^2-w=kx^2.
\tag{Q49.4.4}\label{eq:Q49:4.4}
\]

Put \(v=\allowbreak w-1\ge\allowbreak 0\) and \(t=\allowbreak kx^2\). Then \(w-v=\allowbreak 1\), \(wv=\allowbreak t\), and the two roots \(w,\allowbreak -v\) satisfy \(z^2=\allowbreak z+t\).
Thus
\[
\begin{aligned}
w^{16}+v^{16}
={}&1+16t+104t^2+352t^3+660t^4\\
&+672t^5+336t^6+64t^7+2t^8.
\end{aligned}
\tag{Q49.4.5}\label{eq:Q49:4.5}
\]
For example, the finite identity follows from
\(S_0=\allowbreak 2\), \(S_1=\allowbreak 1\), \(S_n=\allowbreak S_{n-1}+tS_{n-2}\), where
\(S_n=\allowbreak w^n+(-v)^n\).
Since \(v^2\le\allowbreak  wv=\allowbreak t\), we have \(v^{16}\le\allowbreak  t^8\); retaining the first three positive terms gives
\[
w^{16}\ge1+16kx^2+104k^2x^4+352k^3x^6.
\tag{Q49.4.6}\label{eq:Q49:4.6}
\]

For \(0\le\allowbreak  x\le\allowbreak 1/5\), the ratio of consecutive binomial terms after degree five is at most
\(((16-5)/(5+1))x\le\allowbreak 11/30\). Hence
\[
\sum_{j=5}^{16}\binom{16}{j}x^j
\le\frac{4368}{1-11/30}x^5
=\frac{131040}{19}x^5<7000x^5.
\tag{Q49.4.7}\label{eq:Q49:4.7}
\]
Combining \eqref{eq:Q49:4.6} and \eqref{eq:Q49:4.7} yields
\[
w^{16}-1-\alpha[(1+x)^{16}-1-16x]
\ge x^2[A-Bx+Cx^2-Dx^3+Ex^4].
\tag{Q49.4.8}\label{eq:Q49:4.8}
\]
Here \(A,\allowbreak B,\allowbreak C,\allowbreak D,\allowbreak E\) are scalar coefficients:
\[
\begin{aligned}
A&=16k-120\alpha=2187/128>17,\\
B&=560\alpha=17745/64<278,\\
C&=104k^2-1820\alpha=377949/256>1476,\\
D&=7000\alpha=443625/128<3500,\\
E&=352k^3=39397347/1024>38000.
\end{aligned}
\]
Therefore
\[
w^{16}-1-\alpha[(1+x)^{16}-1-16x]
\ge x^2[17-278x+1476x^2-3500x^3+38000x^4].
\tag{Q49.4.9}\label{eq:Q49:4.9}
\]
The polynomial in brackets is
\[
(17-278x+1152x^2)+x^2(324-3500x+38000x^2).
\]
Completing the first square gives
\[
17-278x+1152x^2
=1152(x-139/1152)^2+263/1152\ge263/1152.
\tag{Q49.4.10}\label{eq:Q49:4.10}
\]
The second bracket is positive for every real \(x\), since
\[
3500^2=12250000<49248000=4\cdot324\cdot38000.
\tag{Q49.4.11}\label{eq:Q49:4.11}
\]
Consequently, on the entire interval,
\[
w^{16}-1-\alpha[(1+x)^{16}-1-16x]\ge(263/1152)x^2.
\tag{Q49.4.12}\label{eq:Q49:4.12}
\]

\paragraph{Q49.5 The uniform deficit}
Using the positive square envelope,
\[
\E H^8
\le c^{16}\{w^{16}-(263/1152)x^2\}
=1-(263/13824)c^{14}d.
\tag{Q49.5.1}\label{eq:Q49:5.1}
\]
Since \(111\cdot15=\allowbreak 1665>\allowbreak 1664=\allowbreak 128\cdot13\), we have \(c>\allowbreak 13/15\).
Furthermore,
\[
(13/15)^7=\frac{62748517}{170859375}>\frac4{11},
\qquad
11\cdot62748517=690233687>683437500=4\cdot170859375.
\]
Thus \(c^{14}>\allowbreak 16/121>\allowbreak 1/8\), and
\[
\E H^8\le1-(263/110592)d\le1-d/500,
\tag{Q49.5.2}\label{eq:Q49:5.2}
\]
because \(263\cdot500=\allowbreak 131500>\allowbreak 110592\).
This proves \eqref{eq:Q49:1.3}.
The conclusion allows one-atom and unequal-mass spaces, zero edge factors, and nonzero balanced kernels at \(d=\allowbreak 0\).
At positive-mass atoms, \(p,\allowbreak q>\allowbreak 0\); the argument uses division only by the positive scalar \(c\).
